\section{Properties}
\label{sec-expec-prop}

\begin{enumerate}[label=\arabic*.]
\item \textbf{Expectation is invariant under ``almost-sure
equivalence'':} If $X\le Y$\emph{almost surely}, meaning
$\mathrm{P}(X\le Y)=1$, then by the definition we
have $\mathbf{E} X \le \mathbf{E} Y$, since any
simple function $Z$ such that $Z\le X$ can be
replaced with another simple function $Z'$ such that
$Z'\le Y$ and $Z=Z'$ almost surely. It follows also
that if $X=Y$ almost surely then
$\mathbf{E} X=\mathbf{E} Y$.

\item \textbf{Triangle inequality:}$|\mathbf{E} X| \le \mathbf{E} |X|$.

\begin{proof}

\[
|\mathbf{E} X| = |\mathbf{E} X_+ - \mathbf{E} X_- |  \le \mathbf{E} X_+ + \mathbf{E} X_- = \mathbf{E} |X|.
\]

\end{proof}

\item \textbf{Markov's inequality} (Called Chebyshev's inequality
in [Dur2010]):

\[
\mathrm{P}(X \ge t) \le \frac{\mathbf{E} X}{t}.
\]

\begin{proof}
Use monotonicity twice to deduce:

\[
\mathrm{P}(X \ge t) = \mathbf{E}(1_{\{X\ge t\}}) \le \mathbf{E}\left[ \frac{1}{t} X 1_{\{X\ge t\}} \right] \le \frac{\mathbf{E} X}{t}.
\]

\end{proof}

\item \textbf{Variance:} If $X$ has finite expectation, we define
its $\textbf{variance}$ to be

\[
\mathbf{V}(X) = \mathbf{E}(X-\mathbf{E} X)^2.
\]

\noindent If $\mathbf{V}(X)<\infty$, by expanding the square it
is easy to rewrite the variance as

\[
\mathbf{V}(X) = \mathbf{E}(X^2)-(\mathbf{E} X)^2.
\]

\noindent We denote $\sigma(X) = \sqrt{\mathbf{V}(X)}$ and call
this quantity the $\textbf{standard deviation of } X$.
Note that if $a\in \R$ then
$\mathbf{V}(a X) = a^2 \mathbf{V}(X)$ and
$\sigma(a X)=|a| \sigma(X)$.

\item \textbf{Chebyshev's inequality:}

\[
\mathrm{P} ( | X-\mathbf{E} X| \ge t) \le \frac{\mathbf{V}(X)}{t^2}.
\]

\begin{proof}
Apply Markov's inequality to $Y=(X-\mathbf{E} X)^2$.
\end{proof}

\item \textbf{Cauchy-Schwartz inequality:}

\[
\mathbf{E}|XY| \le \left(\mathbf{E} X^2 \mathbf{E} Y^2\right)^{1/2}.
\]

\noindent Equality holds if and only if $X$ and $Y$ are
linearly dependent, i.e. $aX+bY\equiv 0$ holds for
some $a,b\in \R$.

\begin{proof}
Consider the function

\[
p(t) = \mathbf{E}(|X| + t |Y|)^2 = t^2 \mathbf{E} Y^2 + 2t \mathbf{E}|XY| + \mathbf{E} X^2.
\]

\noindent Since $p(t)=at^2+bt+c$ is a quadratic polynomial in
$t$ that satisfies $p(t)\ge 0$ for all $t$, its
discriminant $b^2-4ac$ must be non-positive. This
gives

\[
(\mathbf{E}|XY|)^2 - \mathbf{E} X^2 \mathbf{E} Y^2 \le 0,
\]

\noindent as claimed. The condition for equality is left as an
exercise.
\end{proof}

\item \textbf{Jensen's inequality:} A function $\varphi:\R\to\R$
is called convex if it satisfies

\[
\varphi(\alpha x + (1-\alpha)y) \le \alpha \varphi(x)+(1-\alpha)\varphi(y)
\]

\noindent for all $x,y\in\R$ and $\alpha\in [0,1]$. If $\varphi$
is convex then

\[
\varphi(\mathbf{E} X) \le \mathbf{E} (\varphi(X)).
\]

\begin{proof}
See homework.
\end{proof}

\item \textbf{$L_p$-norm monotonicity:} If $0<r\le s$ then

\begin{equation*} \tag{6}\label{eq:lpnormmono}
(\mathbf{E}|X|^r)^{1/r} \le (\mathbf{E} |X|^s)^{1/s}.
\end{equation*}

\begin{proof}
Apply Jensen's inequality to the r.v. $|X|^r$ with
the convex function $\varphi(x) = x^{s/r}$.
\end{proof}
\end{enumerate}
